\documentclass[11pt]{article}
\usepackage{amsmath, amssymb}
\usepackage{geometry}
\usepackage{array}
\usepackage{booktabs}
\geometry{margin=0.7in}
\setlength{\parindent}{0pt}
\setlength{\parskip}{0.35em}
\renewcommand{\arraystretch}{1.15}

\newcommand{\cheatsheettitle}[2]{%
\begin{center}
{\LARGE #1}\\[0.4em]
{\large #2}
\end{center}
}


\begin{document}
\cheatsheettitle{Algebra I Linear Functions Cheatsheet}{Module 5}

\section*{Slope}
\[
m=\frac{\text{rise}}{\text{run}}=\frac{y_2-y_1}{x_2-x_1}
\]
\[
\text{positive slope: increasing},\qquad
\text{negative slope: decreasing}
\]
\[
\text{zero slope: horizontal line},\qquad
\text{undefined slope: vertical line}
\]

\section*{Slope-Intercept Form}
\[
y=mx+b
\]
\[
m=\text{slope},\qquad b=\text{\(y\)-intercept}
\]
Graph by plotting $(0,b)$, then using rise over run.

\section*{Point-Slope Form}
\[
y-y_1=m(x-x_1)
\]
Use when you know one point and the slope.

\section*{Standard Form}
\[
Ax+By=C
\]
Intercepts:
\[
x\text{-intercept: set }y=0,\qquad
y\text{-intercept: set }x=0
\]

\section*{Parallel and Perpendicular Lines}
Parallel lines have the same slope:
\[
m_1=m_2
\]
Perpendicular lines have slopes that are opposite reciprocals:
\[
m_1m_2=-1
\]

\section*{Function Notation}
\[
f(x)=mx+b
\]
\[
f(3)\text{ means substitute }x=3
\]

\section*{Writing Linear Equations}
\begin{itemize}
  \item From slope and intercept: use $y=mx+b$.
  \item From slope and point: use $y-y_1=m(x-x_1)$.
  \item From two points: find $m$, then use point-slope form.
  \item From a table: check whether first differences are constant.
\end{itemize}

\section*{Quick Checks}
\begin{itemize}
  \item Slope is change in $y$ over change in $x$.
  \item Vertical lines are $x=a$ and are not functions.
  \item Horizontal lines are $y=b$ and have slope $0$.
  \item The \(y\)-intercept occurs when $x=0$.
\end{itemize}

\end{document}
